%% ma: L11-B04-0017
%% lop: 11
%% chuong: 1
%% bai: 4
%% dang: 11.4.12
%% loai: DS
%% muc: [TH, NB, TH, VD]
%% dap_an: "ĐSĐĐ"
%% nguon: "de-on-tap-giua-ki-1-toan-11-nam-2026-2027-truong-thpt-ngo-quyen-da-nang.pdf (M09, Đề ôn tập giữa kì 1 Toán 11 THPT Ngô Quyền, Đà Nẵng), Mã đề 0111, Phần II Câu 1"
%% kien_thuc: "Bài 2 (công thức lượng giác), Bài 3 (hàm số lượng giác), Bài 4 (phương trình lượng giác cơ bản)"
%% trang_thai: cho-duyet
%% ly_do: "Dạng mới đề xuất 11.4.12"
\begin{ex}
Cho hàm số $f(x)=2\cos\left(2x-\dfrac{\pi}{3}\right)+1$.
\choiceTF{\True Với mọi $x\in\mathbb R$, ta có $2\cos\left(2x-\dfrac{\pi}{3}\right)=\cos2x+\sqrt3\sin2x$}{Giá trị của hàm số tại $x=\dfrac{\pi}{4}$ là $f\left(\dfrac{\pi}{4}\right)=2+\sqrt3$}{\True Tập giá trị của hàm số $f(x)$ là đoạn $[-1;3]$}{\True Phương trình lượng giác $f(x)=2$ có đúng $3$ nghiệm thuộc đoạn $[0;\pi]$}
\loigiai{a) Đúng: $2\cos\left(2x-\dfrac\pi3\right)=2\left(\cos2x\cdot\dfrac12+\sin2x\cdot\dfrac{\sqrt3}2\right)=\cos2x+\sqrt3\sin2x$.

b) Sai: $f\left(\dfrac\pi4\right)=2\cos\dfrac\pi6+1=\sqrt3+1$.

c) Đúng: $-1\le\cos\left(2x-\dfrac\pi3\right)\le1$ nên $-1\le f(x)\le3$, các giá trị đầu mút đều đạt được.

d) Đúng: $\cos\left(2x-\dfrac\pi3\right)=\dfrac12\Leftrightarrow x=\dfrac\pi3+k\pi$ hoặc $x=k\pi$. Trên $[0;\pi]$: $x\in\left\{0;\dfrac\pi3;\pi\right\}$.}
\end{ex}
